%% LyX 2.4.4 created this file.  For more info, see https://www.lyx.org/.
%% Do not edit unless you really know what you are doing.
\documentclass[english]{article}
\usepackage[T1]{fontenc}
\usepackage[latin9]{inputenc}
\usepackage{units}
\usepackage{enumitem}
\usepackage{amsmath}
\usepackage{amsthm}
\usepackage{amssymb}
\usepackage{cancel}
\usepackage{geometry}
\geometry{verbose,tmargin=1in,bmargin=1in,lmargin=1in,rmargin=1in}

\makeatletter
%%%%%%%%%%%%%%%%%%%%%%%%%%%%%% Textclass specific LaTeX commands.
\numberwithin{equation}{section}
\numberwithin{figure}{section}
\newlength{\lyxlabelwidth}      % auxiliary length 

%%%%%%%%%%%%%%%%%%%%%%%%%%%%%% User specified LaTeX commands.
\usepackage{fancyhdr}
\usepackage{lastpage}
\pagestyle{fancy}
\usepackage{enumitem}
\usepackage{ifthen}
\everymath=\expandafter{\the\everymath\displaystyle}
%\DeclareMathSizes{10}{10}{10}{10}
\usepackage{multicol}

\usepackage{tikz}
\usepackage{tikz-3dplot}
\usetikzlibrary{calc,arrows}

\usetikzlibrary{patterns}
\usetikzlibrary{plotmarks}
\usepackage{pgfplots}
\pgfplotsset{compat=newest}
\usepgfplotslibrary{polar}

\usepackage{calc}
\usepackage[nomessages]{fp}

\usepackage{datetime}
\usepackage[super]{nth}
% Added by lyx2lyx
\setlength{\parskip}{\smallskipamount}
\setlength{\parindent}{0pt}

\makeatother

\usepackage{babel}
\begin{document}
\renewcommand{\author}{Dr.\,James Ashe}

\newcommand{\classNum}{331}

\newcommand{\classSec}{A}

\newcommand{\examType}{Exam 4}

\newcommand{\name}{Solutions}

\renewcommand{\date}{\formatdate{22}{11}{2013}}

%%First page's header%%%%%%%%%%%%%%%%%%%%%%
\fancypagestyle{firststyle} %%header settings for first pagr
{
	\fancyhead[L]{MTH \classNum{}(\classSec{}) \author{}\\
	\textbf{\examType{}} ~\date{}}
	\fancyhead[C]{\textbf{\name}}
	%\fancyhead[R]{\textbf{Name:\hspace{5cm}}}
	\setlength{\headsep}{14pt}
}
%%%%%%%%%%%%%%%%%%%%%%%%%%%%%%%%%%%%%%%%%%

%%Next page's header%%%%%%%%%%%%%%%%%%%%%%%
\setlist{nolistsep}
\lhead{MTH \classNum{}(\classSec{})} 
\chead{\textbf{\name}} 
\cfoot{\thepage{}~of~\pageref{LastPage}} 
\setlength{\headsep}{5pt} 
%%%%%%%%%%%%%%%%%%%%%%%%%%%%%%%%%%%%%%%%%%%

%%Various settings; see the preamble for other settings%%%%%
%\setenumerate[0]{leftmargin=12pt,itemindent=0pt,labelwidth=0pt}
\setlist{topsep=.5pt}
\renewcommand{\labelenumi}{\textbf{\arabic{enumi}.}}
\renewcommand{\labelenumii}{\textbf{\alph{enumii}.}}
\thispagestyle{firststyle}
%%%%%%%%%%%%%%%%%%%%%%%%%%%%%%%%%%%%%%%%%%

\newcommand{\xmax}{0}
\newcommand{\xmin}{-\xmax}
\newcommand{\xstep}{0}
\newcommand{\ymax}{0}
\newcommand{\ymin}{-\ymax}
\newcommand{\ystep}{0} 
\newcommand{\dspace}{\vspace{1cm}}
\newcommand{\xa}{0}
\newcommand{\xb}{0}
\newcommand{\ya}{1}
\newcommand{\yb}{1}

\newcommand{\xspace}{\vspace{1cm}}
\newcounter{probs}

\textbf{Unless stated otherwise, each question is of equal value {[}4
pts{]}. Show all of your work.}

\emph{Find the domain and range of the following function.}
\begin{enumerate}[resume]
\item $g(x,y)=\sqrt{x-2y}$

\begin{enumerate}[resume]
\item []Domain: $\left\{ (x,y)\colon x-2y\geq0\right\} $
\item []Range: $[0,\infty)$\addtocounter{probs}{1}\vfill
\end{enumerate}
\end{enumerate}
\emph{Evaluate the following limits.}
\begin{enumerate}[resume]
\item ${\displaystyle \lim_{(x,y)\rightarrow(1,0)}(4x-y})$

\begin{enumerate}[resume]
\item []$=4(1)-0=4$ since the function is continuous on $\mathbb{R}^{2}$.\addtocounter{probs}{1}\vfill
\end{enumerate}
\item ${\displaystyle \lim_{(x,y)\rightarrow(1,-2)}\frac{y^{2}+2xy}{y+2x}}$

\begin{enumerate}[resume]
\item []${\displaystyle \lim_{(x,y)\rightarrow(1,-2)}\frac{y\cancel{y+2x}}{\cancel{y+2x}}=\lim_{(x,y)\rightarrow(1,-2)}y=-2}$\addtocounter{probs}{1}\vfill
\end{enumerate}
\end{enumerate}
\emph{Use the Two-Path Test to prove that the following limit do not
exist.}
\begin{enumerate}[resume]
\item ${\displaystyle \lim_{(x,y)\rightarrow(0,0)}\frac{xy}{x^{2}+y^{2}}}$

\begin{enumerate}[resume]
\item []Take the path $y=x$, then ${\displaystyle \lim_{(x,x)\rightarrow(0,0)}\frac{x^{2}}{2x^{2}}=\frac{1}{2}}$.\\
But if we take the path $y=-x$, then ${\displaystyle \lim_{(x,-x)\rightarrow(0,0)}\frac{-x^{2}}{x^{2}+(-x)^{2}}=\lim_{(x,-x)\rightarrow(0,0)}\frac{-x^{2}}{2x^{2}}=-\frac{1}{2}}$.
So $\lim_{(x,-x)\rightarrow(0,0)}\frac{-x^{2}}{x^{2}+(-x)^{2}}$ DNE.\addtocounter{probs}{1}\vfill
\end{enumerate}
\end{enumerate}
\emph{Find the indicated partial derivative of the following functions.}
\begin{enumerate}[resume]
\item Using $h(x,y)=y^{3}\sin4x$, find $h_{xx}$ and $h_{xy}$.

\begin{enumerate}[resume]
\item []$h_{x}=4y^{3}\cos4x$\\
$h_{xx}=-16y^{3}\sin4x$ and $h_{xy}=12y^{2}\cos4x$.\addtocounter{probs}{1}\vfill
\end{enumerate}
\end{enumerate}
\emph{Use the Chain Rule to find the following derivative. When feasible
express your answer in terms of the independent variable. }
\begin{enumerate}[resume]
\item {[}6 pts{]} $dz/dt$ where $z=x^{3}y-x$, $x=t^{2}$, and $y=2t^{-1}$.

\begin{enumerate}[resume]
\item []$\frac{dz}{dx}=3x^{2}y-1$, $\frac{dx}{dt}=2t$, $\frac{dz}{dy}=x^{3}$
and $\frac{dy}{dt}=-2t^{-2}$.$\frac{dz}{dt}=\frac{dz}{dx}\frac{dx}{dt}+\frac{dz}{dy}\frac{dy}{dt}$\\
So 
\begin{eqnarray*}
\frac{dz}{dt} & = & \left(3x^{2}y-1\right)\left(2t\right)+\left(x^{3}\right)\left(-2t^{-2}\right)\\
 & = & 2t(3t^{4}2t^{-1}-1)-2t^{-2}t^{6}\\
 & = & 12t^{4}-2t-2t^{4}=10t^{4}-2t
\end{eqnarray*}
\addtocounter{probs}{1}
\end{enumerate}
\end{enumerate}
\vfill\newpage

\emph{Use the function }$f(x,y)=x^{2}-4y^{2}-9$\emph{ to complete
problems \ref{enu:Compute-the-gradient}--\ref{enu:Compute-the-directional-END}.}
\begin{enumerate}[resume]
\item \label{enu:Compute-the-gradient}Compute the gradient of $f(x,y)$
at $P(1,-2)$.

\begin{enumerate}[resume]
\item []$f_{x}=2x$ and $f_{y}=-8y$.\\
$\nabla f(1,-2)=\left\langle 2x,-8y\right\rangle |_{(1,-2)}=\left\langle 2,16\right\rangle $\vfill\addtocounter{probs}{1}
\end{enumerate}
\item Find vectors that give the direction of steepest ascent and descent
for $f$ at $P(1,-2)$.

\begin{enumerate}[resume]
\item []ascent: $\nabla f(1,-2)=\left\langle 2,16\right\rangle $ and descent:
$-\nabla f(1,-2)=-\left\langle 2,16\right\rangle $\vfill\addtocounter{probs}{1}
\end{enumerate}
\item \label{enu:Compute-the-directional-END}Compute the directional derivative
of $f(x,y)$ at $P(1,-2)$ in the direction of $\mathbf{u}=\left\langle \frac{1}{\sqrt{5}},\frac{2}{\sqrt{5}}\right\rangle $
(note that $\left|\mathbf{u}\right|=1$).

\begin{enumerate}[resume]
\item []$\left\langle \frac{1}{\sqrt{5}},\frac{2}{\sqrt{5}}\right\rangle \cdot\left\langle 2,16\right\rangle =\frac{2}{\sqrt{5}}+\frac{32}{\sqrt{5}}=\frac{34}{\sqrt{5}}$\vfill\addtocounter{probs}{1}
\end{enumerate}
\end{enumerate}
\vfill\emph{Use the Second Derivative Test to determine (if possible)
whether each critical point of the given function is a local maximum,
local minimum, or saddle point.}
\begin{enumerate}[resume]
\item $f(x,y)=4+x^{3}+y^{3}-3xy$ with critical points: \begin{multicols}{2}

\begin{enumerate}
\item []$f_{x}=3x^{2}-3y$, $f_{xx}=6x$, $f_{y}=3y^{2}-3x$, $f_{yy}=6y$,
and $f_{xy}=-3$\\
$D_{f}(x,y)=36xy-9$
\item $(0,0)$ : $D_{f}(0,0)=0\cdot0-9<0\Rightarrow$ $(0,0)$ is a saddle
point.\vfill
\item $(1,1)$: $D_{f}(1,1)=36-9>0$ and $f_{xx}(1,1)=6>0$ $\Rightarrow$
$(1,1)$ is a local minimum.\vfill
\end{enumerate}
\end{multicols}
\end{enumerate}
\vfill\emph{ Find all of the critical points (there are two) of the
following function.}
\begin{enumerate}[resume]
\item $f(x,y)=\frac{x^{3}}{3}-\frac{y^{3}}{3}+3xy$

\begin{enumerate}[resume]
\item []$f_{x}=x^{2}+3y$ and $f_{y}=-y^{2}+3x$\\
$x^{2}+3y=0\Rightarrow x=0\mbox{ and }y=0\Rightarrow(0,0)$ OR $x^{2}=-3y\Rightarrow y=\nicefrac{-x^{2}}{3}$$\Rightarrow x^{4}-27x=0$
$\Rightarrow x(x^{3}-27)=0$$\Rightarrow x=0$ or $x=3$ thus $(3,-3)$
and $(0,0)$ are critical points.\addtocounter{probs}{1}
\end{enumerate}
\end{enumerate}
\vfill%\arabic{probs}
\end{document}
